\section{Introduction}

Sch\"onhage and Strassen proved in 1971 that two $n$-bit integers can
be multiplied in $O(n\log n\log\log n)$ time on a multitape Turing
machine, and proposed $n\log n$ as the optimal order of growth
\cite[Section~1]{SchonhageStrassen1971}. After a sequence of improvements,
Harvey and van der Hoeven reached the unconditional $O(n\log n)$ upper
bound~\cite{HarveyHoeven2021}. We obtain a strict power saving in the
logarithmic factor, in the same fixed finite-tape model. The algorithm
uses linear combinations of intermediate data to reduce the cost of
both address rearrangement and Fourier-transform layers.

For a binary string $x$ of length $n$, let $\operatorname{val}(x)$ be
its nonnegative integer value, with the leftmost bit most significant.
For $0\le z<2^k$, let $\operatorname{bin}_k(z)$ denote its binary
representation padded on the left to length $k$. We consider exact
multiplication: on input $x\#y$, with $x,y\in\{0,1\}^n$, the output
must be $\operatorname{bin}_{2n}(\operatorname{val}(x)\operatorname{val}(y))$.
Write $\lg n=\max\{\lceil\log_2 n\rceil,1\}$.

\begin{theorem}\label{thm:main}
There is one deterministic Turing machine $A$, with a fixed finite
alphabet and a fixed finite number of one-dimensional tapes, that
computes this exact product for every $n\ge1$ and every pair of $n$-bit
inputs. Its worst-case running time satisfies
\[
 T_A(n)=O\!\left(n(\lg n)^{1-\kappa}\right),
 \qquad \kappa=2^{-182}.
\]
\end{theorem}

The bound is asymptotic through all input lengths. The constants and
thresholds in the construction are extremely large; the algorithm is
intended to establish a bit-complexity bound. The theorem refutes an
eventual $\Omega(n\log n)$ lower bound for every fixed multiplication
machine in this model.

\subsection{The two tape operations}

Splitting the inputs into radix digits reduces multiplication to
convolution of short integer coefficients. Our starting point is the
Gaussian-resampling reduction of Harvey and van der Hoeven
\cite{HarveyHoeven2021}: it replaces the Fourier transforms needed for
this convolution by multidimensional transforms whose axis lengths are
powers of two. After a further conversion to convolution, one coordinate
is stored as the coefficient list of a polynomial modulo $y^r+1$, with $r$ a power of two. Powers of $y$ then
serve as roots of unity in the remaining axes, and multiplication by
such a root is a signed shift of the coefficient list. These are the
synthetic polynomial transforms of Nussbaumer and Quandalle and
Nussbaumer~\cite{NussbaumerQuandalle1978,Nussbaumer1980}.

In this representation, the transform calculation consists of address
rearrangements, signed shifts, and two-point butterfly operations.
Scanning an array of $\Theta(n)$ stored bits at each of
$\Theta(\log n)$ transform levels already costs $\Theta(n\log n)$.
We therefore need savings in both movement and arithmetic. Two tape
procedures supply them; in both bounds, $V$ denotes the stored bit volume.

The first procedure interchanges two disjoint contiguous address fields
of $u$ bits in an array whose addresses form a complete Cartesian
product. Its cost is $O(Vu^\tau)$ for a fixed $\tau<1$.
Although its final effect is an exact permutation, its intermediate
steps form pointwise XOR combinations of data. The final permutation is
exact for arbitrary initial auxiliary values.

The second procedure applies
\[
 H_0(u,v)=\bigl((u+v)/2,(u-v)/2\bigr)
\]
simultaneously on selected coordinate bits of the numerical arrays.
Each record is a polynomial with complex fixed-point coefficients, and
$H_0$ acts coefficientwise. On the layouts used in the multiplication
algorithm, one selects the same bit position in each of at most $d$
equal-width address chunks, and the cost is $O(Vd^{\lambda'})$ for a
fixed $\lambda'<1$. Intermediate values are exact Gaussian dyadics,
whose real and imaginary parts have power-of-two denominators. A single
truncation follows the completed contractive layer.
Section~\ref{sec:tape-interchange} proves the address bound, and
Section~\ref{sec:fast-butterfly-layers} gives the numerical layout,
precision hypotheses, and layer bound.

\subsection{Linear networks and the recursive saving}

Each procedure arises from its own fixed linear network that exchanges two
banks of values, with prescribed signs, and restores every auxiliary
value. A wire's fixed place in the network is called its role. To use
the network on an array, regard a row as the complete ordered suffix,
including all active address fields, at a fixed prefix and row index.
Distribute rows cyclically among the $W$ roles, padding the row count
for each prefix to a multiple of $W$. Let $V$ now denote the node volume
after padding. All role streams have the same local address set and each
contains $V/W$ bits. A scalar gate is then applied
pointwise to matching addresses on the roles that it touches.

The representation of a role's array may change along the network.
At a source, gate, or sink $v$, choose an invertible linear operator
$D_v$ on the space of arrays indexed by the common address set. We call $D_v$ its frame: the stored
array is $D_vg$ when the logical array is $g$. All incidences of one
gate have the same frame. Along an edge from $u$ to $v$, the operator
$D_vD_u^{-1}$ changes the stored array from $D_ug$ to $D_vg$.
At a pointwise linear gate $G$, the common frame commutes with the gate:
\[
 G(D_vg_1,\ldots,D_vg_t)=D_vG(g_1,\ldots,g_t).
\]
The supplied input needs no preliminary change of frame: its logical
interpretation is defined by applying the inverse input frame. Thus the
intermediate frames cancel, and a route from input role $w$ to output
role $\rho(w)$ applies
$D_{\mathrm{sink}(\rho(w))}D_{\mathrm{source}(w)}^{-1}$, together with
its prescribed scalar sign. Undoing the role exchange and those signs
leaves the endpoint operator on every role array, including the roles
that carried auxiliary scalar values.

For the bit procedure, that endpoint operator is an address shear:
two fields, written as vectors $H,D$ over a common residue ring, are
updated by $(H,D)\mapsto(H+D,D)$. Two further linear-cost address
updates complete their interchange. For the numerical procedure, the
endpoint operator becomes the simultaneous $H_0$ layer after diagonal
phase operations and a dyadic scale. The intermediate frames are chosen
so that every edge can be implemented by smaller calls to the same
procedure.

The finite construction in Section~\ref{sec:finite-networks} attaches
subspaces of a fixed $m$-dimensional space to its gates and terminals.
The subspaces at adjacent vertices are nested. Their dimension changes
count the smaller operations needed on an edge, with an additional
source contribution in the bit construction. The construction uses
three-element subsets: gather and scatter create coefficients depending
on intersection size, and auxiliary wires cancel the unwanted
coefficients. Orthogonality of the corresponding indicator vectors permits
the subspace assignments. The decisive property is that the total
number $s$ of smaller calls satisfies $s<Wm$.

A recursive call at parameter $mf$ therefore uses $s$ calls at parameter
$f$, each on a stream of volume $V/W$. The parameter counts field
digits for address interchange and selected axes for a numerical layer.
After dividing cost by volume, multiplying the parameter by $m$
increases the recursive contribution by $s/W<m$. This strict inequality
is the source of the power saving. The later tape and arithmetic
estimates preserve it, including the work of changing numerical address
coordinates. Those changes are implemented by whole-chunk moves and
controlled shifts, followed by a deterministic correction of the
explicitly identified exceptional addresses.

\subsection{From transforms to the exact product}

The remaining proof places these procedures inside the multiplication
reduction and accounts for every change of representation.
First, the radix-digit coefficient list is placed on a product of
distinct prime cyclic axes by a Chinese-remainder map that preserves
convolution. Section~\ref{sec:assembly} implements this map through
triangular controlled shifts and charged axis moves. Gaussian
resampling then expresses each source Fourier transform through a
power-of-two tensor transform. Section~\ref{sec:resampling} retains the
source and target coordinate permutations in the transform identity,
so they need not be executed as separate full sorts.

Next, Bluestein's chirp identity converts the middle transform,
including its target permutation, into normalized cyclic convolution.
A phase twist changes the last cyclic coordinate into a polynomial
modulo $y^r+1$; its negative wrap gives precisely the original cyclic
wrap after the twist is removed. Sections~\ref{sec:synthetic-layouts}
and~\ref{sec:exact-ring-products} compute convolution on the remaining
axes over this polynomial ring. The two forward synthetic transforms
retain the same frequency order through their pointwise product, and
the opposite transform consumes that order. Their polynomial products
use signed packing and the established $O(n\log n)$ multiplier as an
ordinary subroutine.

Finally, the source permutation retained during resampling cancels
across the two forward source transforms, their pointwise product,
and the opposite-sign transform. Section~\ref{sec:assembly} removes
the known normalizations and digit scales, proves that the recovered
coefficient error is strictly below $1/2$, and propagates carries to
obtain the exact binary product. A compatible choice of rational
parameters makes all terms in its cost table fit Theorem~\ref{thm:main}.
Figure~\ref{fig:multiplication-flow} summarizes these reductions.

\begin{figure}[ht]
\centering
\begin{tikzpicture}[
 node distance=10mm,
 every node/.style={font=\small},
 box/.style={draw,rounded corners=2pt,align=center,text width=9.6cm,inner sep=5pt},
 >=Stealth]
\node[box] (digits) {Integer product\\radix-digit convolution};
\node[box,below=of digits] (prime) {Cyclic convolution on distinct prime axes};
\node[box,below=of prime] (power) {Normalized convolution on power-of-two axes};
\node[box,below=of power] (ring) {Convolution over $\mathbb C[y]/(y^r+1)$\\
synthetic transforms and packed polynomial products};
\draw[->] (digits)--node[right,font=\scriptsize,align=left]
 {Chinese-remainder\\address map}(prime);
\draw[->] (prime)--node[right,font=\scriptsize,align=left]
 {Gaussian resampling\\and chirps}(power);
\draw[->] (power)--node[right,font=\scriptsize,align=left]
 {last-coordinate\\phase twist}(ring);
\end{tikzpicture}

\caption{The reductions connecting the tape procedures to multiplication.
The arrows describe changes of problem and representation. Fast address
interchanges organize the layouts, and simultaneous butterfly layers
accelerate the synthetic transforms. Solving the final polynomial-ring
convolution evaluates the transforms needed in the preceding stages.
Known rescalings, rounding, and carry propagation then recover the exact
integer product.}
\label{fig:multiplication-flow}
\end{figure}

\subsection{Other arithmetic operations}

\begin{corollary}[Exact division and integer square root]
\label{cor:exact-arithmetic}
There are deterministic machines with fixed finite alphabets and fixed
finite numbers of one-dimensional tapes performing the following tasks
in $O(n(\lg n)^{1-\kappa})$ worst-case time, for every $n\ge1$. For
$x,y\in\{0,1\}^n$, put $a=\operatorname{val}(x)$ and
$b=\operatorname{val}(y)$. If $b>0$, the division machine takes $x\#y$
and returns $\operatorname{bin}_n(q)\#\operatorname{bin}_n(r)$, where
$q=\lfloor a/b\rfloor$ and $r=a-bq$. The square-root machine takes $x$
and returns $\operatorname{bin}_n(\lfloor\sqrt a\rfloor)$.
Leading zeroes are allowed, and $\kappa$ is as in
Theorem~\ref{thm:main}.
\end{corollary}

The classical Newton reductions described by
Brent~\cite[Lemmas 2.1--2.4]{Brent1976} compute reciprocal and
square-root approximations at geometrically increasing precisions.
Their multiplication costs are bounded by a geometric sum controlled by
the last precision, and exact remainder and square tests recover the
required integer answers. Section~\ref{sec:exact-arithmetic-proof}
gives the complete fixed-tape implementation.


\subsection{Historical context and inherited methods}

Karatsuba's recursive construction, published jointly
with Ofman in 1962, gave a binary multiplication circuit of size
$O(n^{\log_2 3})$~\cite{KaratsubaOfman1962}. Toom's subsequent use of polynomial
evaluation and interpolation gave circuit size
$n\exp(O(\sqrt{\log n}))$, and hence $O(n^{1+\eta})$ for every fixed
$\eta>0$~\cite{Toom1963}. These early results were formulated in terms of
automata and logical networks. Cook adapted Toom's method to the
Turing-machine setting; the contemporary account of Sch\"onhage and
Strassen credits this adaptation explicitly
\cite[Section~1, p.~282]{SchonhageStrassen1971}.

In the multitape Turing model,
Sch\"onhage and Strassen proved the bound
$O(n\log n\log\log n)$ in 1971~\cite{SchonhageStrassen1971}.
F\"urer improved it to $n\log n\,2^{O(\log^* n)}$ in 2007, with a full
journal account in 2009~\cite{Furer2007,Furer2009}; here $\log^* n$ counts
iterated logarithms until a constant threshold is reached.
De, Kurur, Saha and Saptharishi developed a modular-arithmetic route to
the same scale~\cite{DeKururSahaSaptharishi2013}.
Later improvements made the exponential factor $8^{\log^* n}$ and then
$4^{\log^* n}$~\cite{HarveyHoevenLecerf2016,HarveyHoeven2019Lattice}.
Harvey and van der Hoeven proved the unconditional
$O(n\log n)$ bound~\cite{HarveyHoeven2021}.

The multitape bounds just cited concern that fixed model. Sch\"onhage obtained
linear-time multiplication on storage-modification machines~\cite{Schonhage1980}.
Groff later gave a smaller-than-$\log n$ factor in a unit-cost RAM model
with constant-time arbitrary memory access, after polynomial preprocessing
whose cost is excluded from that bound~\cite[Sections~2 and~4]{Groff2019}.
Neither result establishes the bound in the fixed finite-tape model of
Theorem~\ref{thm:main}.

The transform methods behind these bounds also underlie the present
construction. Sch\"onhage--Strassen multiplication uses roots of unity in
Fermat-type residue rings, where certain multiplications become shifts.
F\"urer combines cheap polynomial roots of unity with more general roots,
reducing the frequency of expensive multiplications. Harvey and van der
Hoeven use Gaussian resampling to pass to multidimensional transforms with
power-of-two dimensions, then evaluate them using polynomial transforms
developed by Nussbaumer and Quandalle and subsequently by
Nussbaumer~\cite{NussbaumerQuandalle1978,Nussbaumer1980,HarveyHoeven2021}.
The Gaussian-gridding work of Dutt and Rokhlin is an earlier conceptual
precursor~\cite{DuttRokhlin1993}; Harvey and van der Hoeven discuss the
connection in~\cite[Section~4.4.3]{HarveyHoeven2021}. The quantitative
resampling and bit-cost estimates used here are those of the latter paper.

More specifically, the proof adapts Harvey and van der Hoeven's
Gaussian-resampling identity, contraction bounds, and short-convolution
procedures~\cite[Theorem~4.2, Proposition~4.7 and Lemmas~4.8--4.12]{HarveyHoeven2021}.
Their framework already permits arbitrary dimension. Here the source and
target permutations remain in the transform identity, while the faster tape
procedures perform the required axis movements. Bluestein's conversion to
convolution~\cite{Bluestein1970} and the synthetic-root convolution algorithm
are retained; the latter follows the polynomial-transform construction
in~\cite[Section~2.4 and Section~3]{HarveyHoeven2021}. Signed coefficient
packing is a form of Kronecker substitution, with the fixed-tape version
of~\cite[Lemma~2.5]{HarveyHoeven2021} as a direct antecedent.
Calls to the established $O(n\log n)$ multiplier are ordinary subroutine
calls; they do not invoke the improved theorem recursively.

\subsection{The finite-network ingredients}

The subset-intersection construction in Section~\ref{sec:finite-networks}
has an earlier source outside multiplication. Alon's presentation of the
Frankl--Wilson construction~\cite{FranklWilson1981} gives low-degree
representations of complementary intersection graphs over different fields
\cite[Section~3, remark after the proof of Theorem~1.1]{Alon1998}.
At the prime $p=2$, its subsets have size three and the two polynomial evaluations
are the intersection count modulo two and the intersection count minus one.
These are the two pairing formulas used here. We change the ground-set
size and express the second pairing through a rational bilinear form.

The gather and scatter operations, together with the side wires, have the
form of linear index-code decoding: a low-rank linear summary gives the
wanted coordinate plus contributions from known neighboring coordinates,
and those contributions are subtracted. Bar-Yossef, Birk, Jayram and Kol
establish the fitting-matrix formulation and this decoder
\cite[Section~3, proof of Theorem~5]{BarYossefBirkJayramKol2006}.
Lubetzky and Stav treat decoding over other fields and use related
intersection representations to separate their possible ranks
\cite[Proposition~2.1 and Section~2.3]{LubetzkyStav2009}.
These comparisons explain the algebraic role of the two fields here;
they do not supply the subspace labels on the subsequent network.

Restoring scratch registers with arbitrary initial contents also has a
precise precedent. The transparent-computation construction of Buhrman,
Cleve, Kouck\'y, Loff and Speelman subtracts an initial register value,
performs a computation, adds the resulting value, and reverses the
computation~\cite[Section~3]{BuhrmanCleveKouckyLoffSpeelman2014}.
Our scratch schedule is the linear-readout version of that construction.
The rank reduction used for address shears is a singular-matrix Bruhat
factorization: after reversing the coordinate order, Grigor'ev's
upper-triangular factors become the two lower-triangular factors needed
here~\cite[Appendix~A.2, Proposition~14(c)]{Grigoriev1982}.

The finite networks and their labels are proved directly in
Section~\ref{sec:finite-networks}. The additional requirements are the
three-stage bank exchange, nested subspaces with a strict total rank or
dimension deficit, and endpoint frames acting on every role. The tape
implementations then turn that deficit into the volume-normalized
recurrences of Sections~\ref{sec:tape-interchange}
and~\ref{sec:fast-butterfly-layers}, including the costs of addressing,
padding, coefficient growth, and finite precision. The cited decoding,
restoration, and factorization results provide ingredients; these further
properties are established in the present proof.

\subsection{Lower-bound scope and matrix transposition}

Sch\"onhage and Strassen conjectured $n\log n$ as the optimal order of
multiplication cost~\cite[Section~1]{SchonhageStrassen1971}. In the present
fixed-tape model, an eventual lower bound for every multiplication machine means that, for
each fixed machine $A$, there are constants $c_A>0$ and $N_A$ such that
$T_A(n)\ge c_A n\lg n$ whenever $n\ge N_A$.
Its negation requires one machine satisfying
$\liminf_{n\to\infty}T_A(n)/(n\lg n)=0$.
The power saving proved here gives a limit of zero through all input lengths,
which is a stronger conclusion.

Harvey and van der Hoeven consider a binary transposition machine
supplied an integer $m\ge1$ and a $k\times k$ binary matrix, where
$k=\lfloor\sqrt m\rfloor$. They proved that if every fixed such machine
has worst-case cost $\Omega(m\lg m)$ through all sufficiently large
supplied values of $m$, then every fixed multiplication machine has the
corresponding eventual lower bound~\cite[Theorem~1.2]{HarveyHoeven2025}.
Combining their quantitative reduction~\cite[Corollary~6.2]{HarveyHoeven2025}
with Theorem~\ref{thm:main} gives a deterministic fixed-tape algorithm
that transposes a row-major $k\times k$ binary matrix in
$O(k^2(\lg k)^{1-\kappa})$ time, with $\kappa$ as in
Theorem~\ref{thm:main}. Section~\ref{sec:transposition} proves the more
general rectangular-matrix bound.

This comparison concerns unrestricted computation on a fixed number of
one-dimensional tapes. Our address procedure forms XOR combinations of
intermediate data, even though its final effect is a permutation.
It therefore lies outside models restricted to moving data without such
computations; see~\cite[Section~1.4]{HarveyHoeven2025} for the distinction.

\subsection{Companion application to exact Fourier transforms}

The complex phase network of Proposition~\ref{prop:complex-motif-interface}
is also used in the companion manuscript~\cite[Section~1.1, Theorem~1.1
and Corollary~1.2]{OpenAIExactDFT2026}. That manuscript gives a single
algorithm for the exact discrete Fourier transform at every positive
transform length $n$. As $n\to\infty$, its cost is
$O(n(\log n)^{1-10^{-13}})=o(n\log n)$.

The cost is measured in an exact complex-register model that assigns unit
cost to exact complex field operations. The algorithm accepts every
$x\in\mathbb C^n$, with no bound on its entries or on intermediate complex
magnitudes. Its computation on the input uses only addition, subtraction,
and multiplication by prepared input-independent scalars. The magnitudes
of those scalar coefficients are unrestricted as well.

Integer operations and random access on $O(\log(n+2))$-bit words have unit
cost, with a fixed number of words for each address or integer. The model
supplies one specified Fourier root $\zeta_{D_*}$ of explicitly computable
order $D_*<1024n^3$. Scalar preparation, schedule construction, and index
work are charged. This is a separate application of the finite network,
with its own uniform array implementation. It is neither a bit-complexity
consequence nor a stable floating-point FFT consequence of the fixed-tape
multiplication theorem, Theorem~\ref{thm:main}.

